\makeatletter
\newlabel{ej:entropia-tres}{{1.5}{2}{}{ejemplo.1.5}{}}
\newlabel{ej:entropia-tres@cref}{{[ejemplo][5][]1.5}{[1][2][]2}}
\newlabel{ej:binaria}{{1.6}{2}{Entropía binaria}{ejemplo.1.6}{}}
\newlabel{ej:binaria@cref}{{[ejemplo][6][]1.6}{[1][2][]2}}
\newlabel{fig:binaria}{{1.1}{3}{Función de entropía binaria del \cref {ej:binaria}. Una moneda muy sesgada es casi previsible y contiene poca información por lanzamiento}{figure.1.1}{}}
\newlabel{fig:binaria@cref}{{[figure][1][1]1.1}{[1][2][]3}}
\newlabel{ej:dos-bits}{{1.11}{4}{}{ejemplo.1.11}{}}
\newlabel{ej:dos-bits@cref}{{[ejemplo][11][]1.11}{[1][4][]4}}
\newlabel{eje:1-aditiva}{{1.1}{5}{}{ejercicio.1.1}{}}
\newlabel{eje:1-aditiva@cref}{{[ejercicio][1][1]1.1}{[1][5][]5}}
\newlabel{eje:1-archivos}{{1.2}{5}{}{ejercicio.1.2}{}}
\newlabel{eje:1-archivos@cref}{{[ejercicio][2][1]1.2}{[1][5][]5}}
\newlabel{eje:1-cambridge}{{1.3}{5}{}{ejercicio.1.3}{}}
\newlabel{eje:1-cambridge@cref}{{[ejercicio][3][1]1.3}{[1][5][]5}}
\newlabel{eje:1-binaria}{{1.4}{5}{}{ejercicio.1.4}{}}
\newlabel{eje:1-binaria@cref}{{[ejercicio][4][1]1.4}{[1][5][]5}}
\newlabel{eje:1-geometrica}{{1.5}{5}{}{ejercicio.1.5}{}}
\newlabel{eje:1-geometrica@cref}{{[ejercicio][5][1]1.5}{[1][5][]5}}
\newlabel{eje:1-funcion}{{1.6}{5}{}{ejercicio.1.6}{}}
\newlabel{eje:1-funcion@cref}{{[ejercicio][6][1]1.6}{[1][5][]5}}
\newlabel{ej:asimetria}{{2.5}{7}{Asimetría}{ejemplo.2.5}{}}
\newlabel{ej:asimetria@cref}{{[ejemplo][5][]2.5}{[1][7][]7}}
\newlabel{fig:venn}{{2.1}{9}{Relaciones entre entropías e información mutua. Cada círculo representa la entropía de una variable; la intersección es la información compartida y la unión, la entropía conjunta}{figure.2.1}{}}
\newlabel{fig:venn@cref}{{[figure][1][2]2.1}{[1][8][]9}}
\newlabel{ej:mutua-bits}{{2.15}{9}{}{ejemplo.2.15}{}}
\newlabel{ej:mutua-bits@cref}{{[ejemplo][15][]2.15}{[1][9][]9}}
\newlabel{eje:2-vf}{{2.1}{10}{}{ejercicio.2.1}{}}
\newlabel{eje:2-vf@cref}{{[ejercicio][1][2]2.1}{[1][10][]10}}
\newlabel{eje:2-direcciones}{{2.2}{10}{}{ejercicio.2.2}{}}
\newlabel{eje:2-direcciones@cref}{{[ejercicio][2][2]2.2}{[1][10][]10}}
\newlabel{eje:2-kl-moneda}{{2.3}{10}{}{ejercicio.2.3}{}}
\newlabel{eje:2-kl-moneda@cref}{{[ejercicio][3][2]2.3}{[1][10][]10}}
\newlabel{eje:2-conjunta}{{2.4}{10}{}{ejercicio.2.4}{}}
\newlabel{eje:2-conjunta@cref}{{[ejercicio][4][2]2.4}{[1][10][]10}}
\newlabel{eje:2-galletas}{{2.5}{10}{}{ejercicio.2.5}{}}
\newlabel{eje:2-galletas@cref}{{[ejercicio][5][2]2.5}{[1][10][]10}}
\newlabel{eje:2-gibbs-ln}{{2.6}{11}{}{ejercicio.2.6}{}}
\newlabel{eje:2-gibbs-ln@cref}{{[ejercicio][6][2]2.6}{[1][11][]11}}
\newlabel{eje:2-mutua}{{2.7}{11}{}{ejercicio.2.7}{}}
\newlabel{eje:2-mutua@cref}{{[ejercicio][7][2]2.7}{[1][11][]11}}
\newlabel{ej:clases}{{3.3}{12}{}{ejemplo.3.3}{}}
\newlabel{ej:clases@cref}{{[ejemplo][3][]3.3}{[1][12][]12}}
\newlabel{fig:arbol-kraft}{{3.1}{13}{El código prefijo $C_4$ del \cref {ej:clases} como árbol binario: las palabras son las hojas y ninguna está en el camino hacia otra. La suma de Kraft es $\tfrac 12+\tfrac 14+\tfrac 18+\tfrac 18=1$}{figure.3.1}{}}
\newlabel{fig:arbol-kraft@cref}{{[figure][1][3]3.1}{[1][13][]13}}
\newlabel{ej:huffman}{{3.11}{15}{}{ejemplo.3.11}{}}
\newlabel{ej:huffman@cref}{{[ejemplo][11][]3.11}{[1][15][]15}}
\newlabel{fig:huffman}{{3.2}{15}{Árbol de Huffman del \cref {ej:huffman}. Cada nodo interno lleva la suma de las probabilidades de sus hojas}{figure.3.2}{}}
\newlabel{fig:huffman@cref}{{[figure][2][3]3.2}{[1][15][]15}}
\newlabel{eje:3-dos-huffman}{{3.1}{16}{}{ejercicio.3.1}{}}
\newlabel{eje:3-dos-huffman@cref}{{[ejercicio][1][3]3.1}{[1][16][]16}}
\newlabel{eje:3-diadica}{{3.2}{16}{}{ejercicio.3.2}{}}
\newlabel{eje:3-diadica@cref}{{[ejercicio][2][3]3.2}{[1][16][]16}}
\newlabel{eje:3-huffman5}{{3.3}{17}{}{ejercicio.3.3}{}}
\newlabel{eje:3-huffman5@cref}{{[ejercicio][3][3]3.3}{[1][17][]17}}
\newlabel{eje:3-unica}{{3.4}{17}{}{ejercicio.3.4}{}}
\newlabel{eje:3-unica@cref}{{[ejercicio][4][3]3.4}{[1][17][]17}}
\newlabel{eje:3-shannon}{{3.5}{17}{}{ejercicio.3.5}{}}
\newlabel{eje:3-shannon@cref}{{[ejercicio][5][3]3.5}{[1][17][]17}}
\newlabel{eje:3-bloques}{{3.6}{17}{}{ejercicio.3.6}{}}
\newlabel{eje:3-bloques@cref}{{[ejercicio][6][3]3.6}{[1][17][]17}}
\newlabel{eje:3-no-comprime}{{3.7}{17}{}{ejercicio.3.7}{}}
\newlabel{eje:3-no-comprime@cref}{{[ejercicio][7][3]3.7}{[1][17][]17}}
\newlabel{ej:aab}{{4.3}{19}{}{ejemplo.4.3}{}}
\newlabel{ej:aab@cref}{{[ejemplo][3][]4.3}{[1][18][]19}}
\newlabel{fig:aritmetica}{{4.1}{19}{Codificación aritmética del mensaje «aab» (\cref {ej:aab}). Cada nivel subdivide el intervalo del nivel anterior en proporción a las probabilidades del símbolo siguiente}{figure.4.1}{}}
\newlabel{fig:aritmetica@cref}{{[figure][1][4]4.1}{[1][19][]19}}
\newlabel{ej:aab-palabra}{{4.6}{20}{}{ejemplo.4.6}{}}
\newlabel{ej:aab-palabra@cref}{{[ejemplo][6][]4.6}{[1][19][]20}}
\newlabel{fig:compresor}{{4.2}{21}{Compresión con un modelo de lenguaje. Emisor y receptor usan el mismo modelo; el receptor recupera cada token y lo añade al contexto para obtener la distribución del siguiente}{figure.4.2}{}}
\newlabel{fig:compresor@cref}{{[figure][2][4]4.2}{[1][20][]21}}
\newlabel{ej:bpb}{{4.12}{21}{}{ejemplo.4.12}{}}
\newlabel{ej:bpb@cref}{{[ejemplo][12][]4.12}{[1][21][]21}}
\newlabel{eje:4-bpb-concepto}{{4.1}{22}{}{ejercicio.4.1}{}}
\newlabel{eje:4-bpb-concepto@cref}{{[ejercicio][1][4]4.1}{[1][22][]22}}
\newlabel{eje:4-bac}{{4.2}{22}{}{ejercicio.4.2}{}}
\newlabel{eje:4-bac@cref}{{[ejercicio][2][4]4.2}{[1][22][]22}}
\newlabel{eje:4-descodificar}{{4.3}{22}{}{ejercicio.4.3}{}}
\newlabel{eje:4-descodificar@cref}{{[ejercicio][3][4]4.3}{[1][22][]22}}
\newlabel{eje:4-abbb}{{4.4}{22}{}{ejercicio.4.4}{}}
\newlabel{eje:4-abbb@cref}{{[ejercicio][4][4]4.4}{[1][22][]22}}
\newlabel{eje:4-bpb}{{4.5}{22}{}{ejercicio.4.5}{}}
\newlabel{eje:4-bpb@cref}{{[ejercicio][5][4]4.5}{[1][22][]22}}
\newlabel{eje:4-equivocado}{{4.6}{23}{}{ejercicio.4.6}{}}
\newlabel{eje:4-equivocado@cref}{{[ejercicio][6][4]4.6}{[1][22][]23}}
\newlabel{eje:4-equivocado-teo}{{4.7}{23}{}{ejercicio.4.7}{}}
\newlabel{eje:4-equivocado-teo@cref}{{[ejercicio][7][4]4.7}{[1][23][]23}}
\newlabel{ej:monedas}{{5.3}{25}{}{ejemplo.5.3}{}}
\newlabel{ej:monedas@cref}{{[ejemplo][3][]5.3}{[1][25][]25}}
\newlabel{ej:rlhf}{{5.8}{26}{}{ejemplo.5.8}{}}
\newlabel{ej:rlhf@cref}{{[ejemplo][8][]5.8}{[1][26][]26}}
\newlabel{eje:5-beta}{{5.1}{27}{}{ejercicio.5.1}{}}
\newlabel{eje:5-beta@cref}{{[ejercicio][1][5]5.1}{[1][27][]27}}
\newlabel{eje:5-monedas}{{5.2}{27}{}{ejercicio.5.2}{}}
\newlabel{eje:5-monedas@cref}{{[ejercicio][2][5]5.2}{[1][27][]27}}
\newlabel{eje:5-landauer}{{5.3}{28}{}{ejercicio.5.3}{}}
\newlabel{eje:5-landauer@cref}{{[ejercicio][3][5]5.3}{[1][27][]28}}
\newlabel{eje:5-rlhf}{{5.4}{28}{}{ejercicio.5.4}{}}
\newlabel{eje:5-rlhf@cref}{{[ejercicio][4][5]5.4}{[1][28][]28}}
\newlabel{eje:5-lagrange}{{5.5}{28}{}{ejercicio.5.5}{}}
\newlabel{eje:5-lagrange@cref}{{[ejercicio][5][5]5.5}{[1][28][]28}}
\newlabel{eje:5-dpo}{{5.6}{28}{}{ejercicio.5.6}{}}
\newlabel{eje:5-dpo@cref}{{[ejercicio][6][5]5.6}{[1][28][]28}}
\newlabel{ap:soluciones}{{A}{29}{Soluciones de los ejercicios}{appendix.A}{}}
\newlabel{ap:soluciones@cref}{{[appendix][1][2147483647]A}{[1][29][]29}}
\newlabel{sol:eje:1-aditiva}{{A.1}{29}{Tema 1}{section*.4}{}}
\newlabel{sol:eje:1-aditiva@cref}{{[subappendix][1][2147483647,1]A.1}{[1][29][]29}}
\newlabel{sol:eje:1-archivos}{{A.1}{29}{Tema 1}{section*.5}{}}
\newlabel{sol:eje:1-archivos@cref}{{[subappendix][1][2147483647,1]A.1}{[1][29][]29}}
\newlabel{sol:eje:1-cambridge}{{A.1}{30}{Tema 1}{section*.6}{}}
\newlabel{sol:eje:1-cambridge@cref}{{[subappendix][1][2147483647,1]A.1}{[1][30][]30}}
\newlabel{sol:eje:1-binaria}{{A.1}{30}{Tema 1}{section*.7}{}}
\newlabel{sol:eje:1-binaria@cref}{{[subappendix][1][2147483647,1]A.1}{[1][30][]30}}
\newlabel{sol:eje:1-geometrica}{{A.1}{30}{Tema 1}{section*.8}{}}
\newlabel{sol:eje:1-geometrica@cref}{{[subappendix][1][2147483647,1]A.1}{[1][30][]30}}
\newlabel{sol:eje:1-funcion}{{A.1}{31}{Tema 1}{section*.9}{}}
\newlabel{sol:eje:1-funcion@cref}{{[subappendix][1][2147483647,1]A.1}{[1][31][]31}}
\newlabel{sol:eje:2-vf}{{A.2}{31}{Tema 2}{section*.10}{}}
\newlabel{sol:eje:2-vf@cref}{{[subappendix][2][2147483647,1]A.2}{[1][31][]31}}
\newlabel{sol:eje:2-direcciones}{{A.2}{32}{Tema 2}{section*.11}{}}
\newlabel{sol:eje:2-direcciones@cref}{{[subappendix][2][2147483647,1]A.2}{[1][32][]32}}
\newlabel{sol:eje:2-kl-moneda}{{A.2}{32}{Tema 2}{section*.12}{}}
\newlabel{sol:eje:2-kl-moneda@cref}{{[subappendix][2][2147483647,1]A.2}{[1][32][]32}}
\newlabel{sol:eje:2-conjunta}{{A.2}{33}{Tema 2}{section*.13}{}}
\newlabel{sol:eje:2-conjunta@cref}{{[subappendix][2][2147483647,1]A.2}{[1][33][]33}}
\newlabel{sol:eje:2-galletas}{{A.2}{33}{Tema 2}{section*.14}{}}
\newlabel{sol:eje:2-galletas@cref}{{[subappendix][2][2147483647,1]A.2}{[1][33][]33}}
\newlabel{sol:eje:2-gibbs-ln}{{A.2}{34}{Tema 2}{section*.15}{}}
\newlabel{sol:eje:2-gibbs-ln@cref}{{[subappendix][2][2147483647,1]A.2}{[1][34][]34}}
\newlabel{sol:eje:2-mutua}{{A.2}{34}{Tema 2}{section*.16}{}}
\newlabel{sol:eje:2-mutua@cref}{{[subappendix][2][2147483647,1]A.2}{[1][34][]34}}
\newlabel{sol:eje:3-dos-huffman}{{A.3}{35}{Tema 3}{section*.17}{}}
\newlabel{sol:eje:3-dos-huffman@cref}{{[subappendix][3][2147483647,1]A.3}{[1][35][]35}}
\newlabel{sol:eje:3-diadica}{{A.3}{35}{Tema 3}{section*.18}{}}
\newlabel{sol:eje:3-diadica@cref}{{[subappendix][3][2147483647,1]A.3}{[1][35][]35}}
\newlabel{sol:eje:3-huffman5}{{A.3}{36}{Tema 3}{section*.19}{}}
\newlabel{sol:eje:3-huffman5@cref}{{[subappendix][3][2147483647,1]A.3}{[1][36][]36}}
\newlabel{sol:eje:3-unica}{{A.3}{36}{Tema 3}{section*.20}{}}
\newlabel{sol:eje:3-unica@cref}{{[subappendix][3][2147483647,1]A.3}{[1][36][]36}}
\newlabel{sol:eje:3-shannon}{{A.3}{37}{Tema 3}{section*.21}{}}
\newlabel{sol:eje:3-shannon@cref}{{[subappendix][3][2147483647,1]A.3}{[1][36][]37}}
\newlabel{sol:eje:3-bloques}{{A.3}{37}{Tema 3}{section*.22}{}}
\newlabel{sol:eje:3-bloques@cref}{{[subappendix][3][2147483647,1]A.3}{[1][37][]37}}
\newlabel{sol:eje:3-no-comprime}{{A.3}{37}{Tema 3}{section*.23}{}}
\newlabel{sol:eje:3-no-comprime@cref}{{[subappendix][3][2147483647,1]A.3}{[1][37][]37}}
\newlabel{sol:eje:4-bpb-concepto}{{A.4}{38}{Tema 4}{section*.24}{}}
\newlabel{sol:eje:4-bpb-concepto@cref}{{[subappendix][4][2147483647,1]A.4}{[1][38][]38}}
\newlabel{sol:eje:4-bac}{{A.4}{38}{Tema 4}{section*.25}{}}
\newlabel{sol:eje:4-bac@cref}{{[subappendix][4][2147483647,1]A.4}{[1][38][]38}}
\newlabel{sol:eje:4-descodificar}{{A.4}{39}{Tema 4}{section*.26}{}}
\newlabel{sol:eje:4-descodificar@cref}{{[subappendix][4][2147483647,1]A.4}{[1][39][]39}}
\newlabel{sol:eje:4-abbb}{{A.4}{39}{Tema 4}{section*.27}{}}
\newlabel{sol:eje:4-abbb@cref}{{[subappendix][4][2147483647,1]A.4}{[1][39][]39}}
\newlabel{sol:eje:4-bpb}{{A.4}{40}{Tema 4}{section*.28}{}}
\newlabel{sol:eje:4-bpb@cref}{{[subappendix][4][2147483647,1]A.4}{[1][40][]40}}
\newlabel{sol:eje:4-equivocado}{{A.4}{40}{Tema 4}{section*.29}{}}
\newlabel{sol:eje:4-equivocado@cref}{{[subappendix][4][2147483647,1]A.4}{[1][40][]40}}
\newlabel{sol:eje:4-equivocado-teo}{{A.4}{40}{Tema 4}{section*.30}{}}
\newlabel{sol:eje:4-equivocado-teo@cref}{{[subappendix][4][2147483647,1]A.4}{[1][40][]40}}
\newlabel{sol:eje:5-beta}{{A.5}{41}{Tema 5}{section*.31}{}}
\newlabel{sol:eje:5-beta@cref}{{[subappendix][5][2147483647,1]A.5}{[1][41][]41}}
\newlabel{sol:eje:5-monedas}{{A.5}{41}{Tema 5}{section*.32}{}}
\newlabel{sol:eje:5-monedas@cref}{{[subappendix][5][2147483647,1]A.5}{[1][41][]41}}
\newlabel{sol:eje:5-landauer}{{A.5}{42}{Tema 5}{section*.33}{}}
\newlabel{sol:eje:5-landauer@cref}{{[subappendix][5][2147483647,1]A.5}{[1][42][]42}}
\newlabel{sol:eje:5-rlhf}{{A.5}{42}{Tema 5}{section*.34}{}}
\newlabel{sol:eje:5-rlhf@cref}{{[subappendix][5][2147483647,1]A.5}{[1][42][]42}}
\newlabel{sol:eje:5-lagrange}{{A.5}{42}{Tema 5}{section*.35}{}}
\newlabel{sol:eje:5-lagrange@cref}{{[subappendix][5][2147483647,1]A.5}{[1][42][]42}}
\newlabel{sol:eje:5-dpo}{{A.5}{43}{Tema 5}{section*.36}{}}
\newlabel{sol:eje:5-dpo@cref}{{[subappendix][5][2147483647,1]A.5}{[1][43][]43}}
\makeatother
